(1) 颈内流体近似整体运动，其质量 $M_n=\rho A\ell_e$。取 $\xi>0$ 为流体向腔内推进，腔体体积净减小 $\Delta V=-A\xi$，绝热体积压缩导致腔内压强增量 $\Delta p=\rho c^2 A\xi/V$，因 $c^2=(\partial p/\partial\rho)_S$，故颈流体受向外恢复力 $-A\Delta p$。
(2) 牛顿方程 $M_n\ddot\xi+(\rho c^2 A^2/V)\xi=0$，因此 $\omega_H^2=(\rho c^2 A^2/V)/(\rho A\ell_e)=c^2A/(V\ell_e)$。动能 $M_n\dot\xi^2/2$，腔体压缩势能 $(\rho c^2 A^2/V)\xi^2/2$，合起来 $E=M_n\dot\xi^2/2+M_n\omega_H^2\xi^2/2$ 恒定。
(3) 加阻尼后方程为 $M_n\ddot\xi+b\dot\xi+M_n\omega_H^2\xi=0$。当 $b<2M_n\omega_H$ 时解的衰减包络为 $e^{-bt/(2M_n)}$，阻尼频率 $\omega_d=\sqrt{\omega_H^2-b^2/(4M_n^2)}$；弱阻尼品质因数 $Q\simeq M_n\omega_H/b$，其正值反映能量缓慢耗散。
(4) 由 $\omega_H=c\sqrt{A/(V\ell_e)}$，增大 $V$ 或 $\ell_e$ 会降低频率，增大 $A$ 则升高频率。有效长度需包括颈口端部修正；若波长不再远大于腔和颈的尺寸，压力不再均匀、颈流体也不再刚体运动，必须用分布波动理论替代。
